% ===== 原文 pp.104--105（tr-08b：致谢 + 缩写表）=====
% [原文 p.104]
\section*{致谢}
\addcontentsline{toc}{section}{致谢}

首先，特别要感谢 Jaume Gomis：十年前他引导我走入了 AGT 对应的基础；也要感谢我在 AGT 相关论文上的其他合作者，多年来与我分享洞见：Alexander Gorsky、Alexey Milekhin、Yiwen Pan、Wolfger Peelaers、Nikita Sopenko 与 Gustavo Turiaci。我要感谢 Antoine Bourget、Ioana Coman、Taro Kimura、Fabrizio Nieri、Wolfger Peelaers、Shlomo Razamat 与 Jaewon Song，他们热心地提供了文献或评论；还要感谢 Elli Pomoni 与 Alessandro Sfondrini 组织 YRISW 2020。最后，我感谢我的妻子和孩子们，在我准备并后来修订这些讲义期间所给予的耐心。

% [原文 p.105]
\section*{缩写表}
\addcontentsline{toc}{section}{缩写表}

% [Tutor Note: 原文 p.105 的缩写表是"两栏并排、按字母序往下走"的版式（左栏 ABJM--GW，右栏 IR--YRISW，共 33 条）。
%  此处改用两列 longtable：缩写列 + "英文全称（中文释义）"列，字母序不变、条目一条不少，只为中文释义所需宽度。
%  两个最易混的条目见下面的批注。]

{\small
\begin{longtable}{@{}l@{\hspace{1.2em}}p{13.1cm}@{}}
\textbf{ABJM} & Aharony--Bergman--Jafferis--Maldacena（M2 膜世界体积理论）\\
\textbf{AD} & Argyres--Douglas（强耦合的 4d $\cN=2$ 理论）\\
\textbf{ADHM} & Atiyah--Drinfeld--Hitchin--Manin（瞬子的构造）\\
\textbf{AGT} & Alday--Gaiotto--Tachikawa（4d/2d 对应）\\
\textbf{ALE} & asymptotically locally Euclidean space（渐近局部欧氏空间；$\C^2/\Gamma$ 的消解）\\
\textbf{BMT} & Bonelli--Maruyoshi--Tanzini（2d CFT 中的非正规态）\\
\textbf{BPS} & Bogomol'nyi--Prasad--Sommerfield（超对称）\\
\textbf{CFT} & conformal field theory（共形场论）\\
\textbf{COHA} & cohomological Hall algebra（上同调 Hall 代数）\\
\textbf{DAHA} & double affine Hecke algebra（双仿射 Hecke 代数）\\
\textbf{DIM} & Ding--Iohara--Miki（代数）\\
\textbf{DOZZ} & Dorn--Otto--Zamolodchikov--Zamolodchikov（Liouville CFT 的三点函数）\\
\textbf{DS} & Drinfeld--Sokolov（W 代数的约化）\\
\textbf{FI} & Fayet--Iliopoulos（超对称作用量中的参数）\\
\textbf{GW} & Gukov--Witten（面缺陷）\\
\textbf{IR} & infra-red（红外；低能/长距离）\\
\textbf{JK} & Jeffrey--Kirwan（留数规定）\\
\textbf{KK} & Kaluza--Klein（圆上的约化）\\
\textbf{LMNS} & Losev--Moore--Nekrasov--Shatashvili（瞬子配分函数的公式）\\
\textbf{NS} & Nekrasov--Shatashvili（$b\to0$ 极限，即 $\epsilon_1\to0$）\\
\textbf{OPE} & operator product expansion（算子乘积展开）\\
\textbf{QFT} & quantum field theory（量子场论）\\
\textbf{RG} & renormalization group（重整化群）\\
\textbf{SCFT} & superconformal field theory（超共形场论）\\
\textbf{SQCD} & super-QCD（SYM 加上物质）\\
\textbf{SW} & Seiberg--Witten（给出 4d $\cN=2$ 理论红外描述与预势 $F$ 的曲线 $\Sigma$ 与微分 $\lambda$）\\
\textbf{SYM} & super-Yang--Mills（本文指 4d $\cN=2$ 矢量多重态）\\
\textbf{TQFT} & topological quantum field theory（拓扑量子场论）\\
\textbf{UV} & ultra-violet（紫外；高能/短距离）\\
\textbf{VEV} & vacuum expectation value（真空期望值）\\
\textbf{YBE} & Yang--Baxter equation（Yang--Baxter 方程）\\
\textbf{YM} & Yang--Mills（非超对称）\\
\textbf{YRISW} & Young Researchers Integrability School and Workshop（青年研究者可积性学校暨研讨会）\\
\end{longtable}
}

% [Tutor Note: p.105 缩写表里最容易混的两组：(i) DOZZ（Dorn--Otto--Zamolodchikov--Zamolodchikov，Liouville 三点函数）与 Dotsenko--Fateev（DF，自由场/库仑气体表示）不是一回事------p.102 讨论矩阵模型时用的是 DF 表示，而三点函数的最终表达式与归一化才叫 DOZZ；
%  (ii) LMNS（Losev--Moore--Nekrasov--Shatashvili，NS 极限下瞬子配分函数的显式公式）与 NS（Nekrasov--Shatashvili 极限 $\epsilon_2\to0$）虽出自同一组人，但一个是"公式"、一个是"极限"，读正文时须区分。]

% [Tutor Note: 条目 SW 的 $\Sigma,\lambda$ 是 Seiberg--Witten 数据：$\Sigma$ 为 SW 曲线、$\lambda$ 为 SW 微分，$\lambda$ 的周期给出预势导数（$a=\oint\lambda$、$a_D=\oint_{\Sigma_D}\lambda$）；注意本译本的约定 $\epsilon_1=b/r$、$\epsilon_2=1/(rb)$（原文 p.8 与 p.52），与附录 A 的 $b$ 是同一个 $b$。
%  "JK"（Jeffrey--Kirwan 留数规定）与 "LMNS" 分别解决不同问题：JK 给的是从围道积分里"取哪些留数"的规定，LMNS 给的是 NS 极限下配分函数的闭式，两者经常一起出现但不可互换。]
